% Generated draft: requires mathematical and source review.
\documentclass[11pt]{article}
\usepackage[margin=25mm]{geometry}
\usepackage{fontspec}
\setmainfont{Latin Modern Roman}
\usepackage{amsmath,amssymb,mathtools,unicode-math}
\setmathfont{Latin Modern Math}
\usepackage{enumitem,needspace,xurl,hyperref}
\setlength{\parindent}{0pt}
\setlength{\parskip}{5pt}
\setlength{\emergencystretch}{3em}
\newcommand{\N}{\mathbb{N}}
\newcommand{\Z}{\mathbb{Z}}
\newcommand{\Q}{\mathbb{Q}}
\newcommand{\R}{\mathbb{R}}
\newcommand{\C}{\mathbb{C}}
\newcommand{\F}{\mathbb{F}}
\newcommand{\A}{\mathbb{A}}
\newcommand{\PP}{\mathbb{P}}
\newcommand{\E}{\mathbb{E}}
\newcommand{\eps}{\varepsilon}
\newcommand{\dd}{\,\mathrm{d}}
\begin{document}
\section*{UnsolvedMath ID: 800618}
\textbf{Source problem code:} catalogue entry
\subsection{Definitions and mathematical statement}
For infinite subsets \(A,B\subseteq\N\), write \(A\subseteq^{*}B\) if \(A\setminus B\) is finite. An infinite set \(A\subseteq\N\) is a pseudointersection of a family \(\mathcal F\subseteq[\N]^{\aleph_0}\) if \(A\subseteq^{*}F\) for every \(F\in\mathcal F\). The family \(\mathcal F\) has the strong finite intersection property if the intersection of each finite subfamily is infinite. Define
\[
\mathfrak p=\min\bigl\{|\mathcal F|:\mathcal F\subseteq[\N]^{\aleph_0}\text{ has the strong finite intersection property and has no infinite pseudointersection}\bigr\}.
\]
A tower is a sequence \(\langle T_\alpha:\alpha<\kappa\rangle\) of infinite subsets of \(\N\) such that, whenever \(\alpha<\beta<\kappa\),
\[
T_\beta\subseteq^{*}T_\alpha
\quad\text{and}\quad
T_\alpha\not\subseteq^{*}T_\beta,
\]
and the family \(\{T_\alpha:\alpha<\kappa\}\) has no infinite pseudointersection. Define \(\mathfrak t\) to be the least cardinal \(\kappa\) for which such a tower exists.

The supplied source asks the Minimal Tower Problem:
\[
\text{Is }\mathsf{ZFC}+\bigl(\mathfrak p<\mathfrak t\bigr)\text{ consistent?}
\]
Here \(\mathsf{ZFC}\) is the conventional background theory adopted for this draft; the supplied excerpt does not explicitly name a background theory or define the two cardinal characteristics. These reconstructed definitions and this interpretation require review. The source describes this as the main open problem in the field and says that it motivated the study of \(\tau\)-covers. Those descriptions, and the registry's supplied open label, are source qualifications rather than an independently verified current mathematical assessment.
\subsection{Short English statement}
The source asks whether it is consistent that the smallest size of a family of infinite subsets of \(\N\) with the strong finite intersection property and no infinite pseudointersection is strictly smaller than the smallest size of a tower. This draft interprets consistency relative to \(\mathsf{ZFC}\); the excerpt leaves that background implicit.
\subsection{Sources}
\begin{itemize}
\item Local UnsolvedMath registry, permanent ID \(800618\), exact source code \texttt{catalogue entry}; supplied title: ``Selection Principles in Mathematics --- Problem 46 ([vD]).'' The supplied record labels the problem open and was updated on 2026-10-16; this label is not independently verified here.
\item Tsaban, \emph{Selection Principles and Special Sets of Reals: Open Problems} (2006), Problem 46, as identified in the supplied source record: \url{https://arxiv.org/abs/math/0609601}. The record reports access on 2026-10-16 and extraction by \texttt{arxiv\_latex}; these are supplied provenance details, not a claim of independent access or verification.
\item The registry lists Tsaban as proposer, with proposed year 2006, and gives the attribution/context marker \texttt{[vD]}. No expansion of that marker or separate bibliographic reference identifying its referent is supplied.
\end{itemize}
\end{document}
